\documentclass[12pt]{article}
\usepackage{amsmath, amssymb}
\usepackage[margin=1in]{geometry}
\setlength{\parskip}{6pt}
\title{QUBO Formulation: Exact Cover by 3-Sets Problem}
\date{}
\begin{document}
\maketitle

\subsection*{Exact Cover by 3-Sets Problem}

\paragraph{Problem Statement.}
Let $U$ be a finite universe of $N$ elements, and let $\mathcal{S} = \{S_0, S_1, \dots, S_{M-1}\}$ be a collection of $M$ subsets of $U$, where each $S_j \subseteq U$ satisfies $|S_j| = 3$. The Exact Cover by 3-Sets Problem asks whether there exists a subcollection $\mathcal{C} \subseteq \mathcal{S}$ such that every element $i \in U$ belongs to exactly one set in $\mathcal{C}$. If such a subcollection exists, it is called an exact cover. The optimization task is to select subsets to minimize the total coverage violations across all universe elements, where a violation count of zero indicates a valid exact cover.

\paragraph{Decision Variables.}
For each set $S_j$ with index $j \in \{0, \dots, M-1\}$, introduce a binary decision variable $x_j \in \{0,1\}$. The variable $x_j = 1$ indicates that set $S_j$ is selected for the cover, and $x_j = 0$ indicates it is not. The decision vector is denoted $\mathbf{x} = (x_0, x_1, \dots, x_{M-1}) \in \{0,1\}^M$.

\paragraph{QUBO Derivation.}
\textbf{Step 1 --- Original Objective.} The objective is to minimize the sum of squared deviations from the exact cover condition across all universe elements. We define the Hamiltonian $H(\mathbf{x})$ directly as this sum:
\begin{align}
H(\mathbf{x}) = \sum_{i \in U} \left( \sum_{j: i \in S_j} x_j - 1 \right)^2.
\end{align}

\textbf{Step 2 --- Constraint Form.} For each universe element $i \in U$, the exact cover requirement is enforced by the linear equality:
\begin{align}
\sum_{j: i \in S_j} x_j = 1, \quad \forall i \in U.
\end{align}
Let $\mathcal{N}_i = \{ j \in \{0, \dots, M-1\} : i \in S_j \}$ denote the set of indices of subsets containing element $i$. The constraint is equivalently written as $\sum_{j \in \mathcal{N}_i} x_j = 1$.

\textbf{Step 3 --- Penalty Expansion.} We expand the squared penalty for each element $i \in U$ using the idempotent property $x_j^2 = x_j$ for binary variables:
\begin{align}
\left( \sum_{j \in \mathcal{N}_i} x_j - 1 \right)^2 &= \left( \sum_{j \in \mathcal{N}_i} x_j \right)^2 - 2 \sum_{j \in \mathcal{N}_i} x_j + 1 \nonumber \\
&= \sum_{j \in \mathcal{N}_i} x_j^2 + \sum_{\substack{j,k \in \mathcal{N}_i \\ j \neq k}} x_j x_k - 2 \sum_{j \in \mathcal{N}_i} x_j + 1 \nonumber \\
&= \sum_{j \in \mathcal{N}_i} x_j + \sum_{\substack{j,k \in \mathcal{N}_i \\ j \neq k}} x_j x_k - 2 \sum_{j \in \mathcal{N}_i} x_j + 1 \nonumber \\
&= 1 - \sum_{j \in \mathcal{N}_i} x_j + \sum_{\substack{j,k \in \mathcal{N}_i \\ j \neq k}} x_j x_k.
\end{align}

\textbf{Step 4 --- Combined Hamiltonian.} Summing the expanded penalties over all $i \in U$ yields:
\begin{align}
H(\mathbf{x}) &= \sum_{i \in U} \left( 1 - \sum_{j \in \mathcal{N}_i} x_j + \sum_{\substack{j,k \in \mathcal{N}_i \\ j \neq k}} x_j x_k \right) \nonumber \\
&= \sum_{i \in U} 1 - \sum_{j=0}^{M-1} \left( \sum_{i \in U} \mathbb{I}(i \in S_j) \right) x_j + \sum_{j=0}^{M-1} \sum_{\substack{k=0 \\ k \neq j}}^{M-1} \left( \sum_{i \in U} \mathbb{I}(i \in S_j \cap S_k) \right) x_j x_k,
\end{align}
where $\mathbb{I}(\cdot)$ is the indicator function. Since each set has exactly three elements, $\sum_{i \in U} \mathbb{I}(i \in S_j) = 3$. The quadratic term sums over ordered pairs $(j,k)$ with $j \neq k$. Converting to the standard upper-triangular QUBO form $\sum_{j < k}$ introduces a factor of $2$ for each unordered pair:
\begin{align}
H(\mathbf{x}) = N - 3 \sum_{j=0}^{M-1} x_j + \sum_{j=0}^{M-1} \sum_{\substack{k=0 \\ k > j}}^{M-1} 2 |S_j \cap S_k| x_j x_k.
\end{align}

\textbf{Step 5 --- Q Matrix and Offset.} Reading off the coefficients from the general closed form above, the QUBO matrix entries $Q_{j,k}$ and constant offset $c$ are:
\begin{align}
Q_{j,j} &= -3, \quad \forall j \in \{0, \dots, M-1\}, \\
Q_{j,k} &= 2 |S_j \cap S_k|, \quad \forall j < k, \\
c &= N.
\end{align}

\paragraph{Penalty Weight Justification.}
The problem specification defines the objective function itself as the sum of squared penalties for the exact cover constraints. Consequently, no external penalty weight $\lambda$ is introduced or required. The quadratic form inherently penalizes any deviation from the constraint $\sum_{j \in \mathcal{N}_i} x_j = 1$ by exactly the square of the violation magnitude. Because the objective is to minimize this sum, any infeasible assignment incurs a strictly positive energy cost proportional to the squared violation, ensuring that feasible exact covers (energy $0$) are strictly preferred over infeasible ones without the need for heuristic weight tuning.

\paragraph{Correctness Argument.}
Minimizing $H(\mathbf{x})$ over $\mathbf{x} \in \{0,1\}^M$ recovers the optimal solution because $H(\mathbf{x})$ is constructed as a sum of non-negative squared terms. For any feasible exact cover, each constraint is satisfied exactly, yielding $H(\mathbf{x}) = 0$. Since $H(\mathbf{x}) \geq 0$ for all $\mathbf{x}$, the global minimum must be $0$. Any infeasible assignment violates at least one constraint, resulting in $H(\mathbf{x}) > 0$. Thus, the ground state of $H$ corresponds exactly to a valid exact cover if one exists, and the minimum energy value directly indicates feasibility (zero) or the total squared violation otherwise.

\paragraph{Illustrative Example.}
Consider a small instance with $N=4$ universe elements $U=\{0,1,2,3\}$ and $M=3$ sets $S_0=\{0,1,2\}$, $S_1=\{0,1,3\}$, $S_2=\{0,2,3\}$. Each set has size $3$. The pairwise intersections are $|S_0 \cap S_1| = 2$, $|S_0 \cap S_2| = 2$, and $|S_1 \cap S_2| = 2$. Instantiating the general formulas yields:
\begin{align}
Q_{0,0} = Q_{1,1} = Q_{2,2} = -3, \quad Q_{0,1} = Q_{0,2} = Q_{1,2} = 4, \quad c = 4.
\end{align}
The concrete Hamiltonian is:
\begin{align}
H(\mathbf{x}) = 4 - 3(x_0 + x_1 + x_2) + 4(x_0 x_1 + x_0 x_2 + x_1 x_2).
\end{align}
Evaluating $H(\mathbf{x})$ over all $2^3 = 8$ binary assignments gives:
\begin{align}
\begin{array}{c|c}
\mathbf{x} & H(\mathbf{x}) \\ \hline
(0,0,0) & 4 \\
(1,0,0), (0,1,0), (0,0,1) & 1 \\
(1,1,0), (1,0,1), (0,1,1) & 2 \\
(1,1,1) & 1
\end{array}
\end{align}
The minimum energy is $1$, achieved by $\mathbf{x}=(1,0,0)$, $\mathbf{x}=(0,1,0)$, and $\mathbf{x}=(0,0,1)$. These correspond to selecting exactly one set, which covers three elements exactly once and leaves one element uncovered, resulting in a total squared violation of $1$. No assignment yields energy $0$, confirming that no exact cover exists for this instance.

\end{document}